% Paste these three frames immediately after "Our Modified Segregation Metric".
% Uses only packages and colors already present in your supplied preamble.
% Hypothetical populations, not empirical data. Each bar reports a neighborhood
% proportion; positions within bars do not represent household locations.

\begin{frame}{Two Cities with Different Population Shares}
\small
Each city has 10 neighborhoods with 100 residents each.
We study the same racial group $r$ in both cities.
\vspace{0.2cm}

\begin{columns}[T]
\begin{column}{0.49\textwidth}
\centering
\textbf{City H: large group, evenly spread}\\[3pt]
800 of 1,000 residents belong to $r$.
\par\vspace{0.15cm}
\begin{tikzpicture}[x=0.88cm,y=0.75cm]
  \foreach \x in {0,...,4} {
    \foreach \y in {0,1} {
      \fill[gray!20] (\x,\y) rectangle +(0.9,0.82);
      \fill[MyTeal] (\x,\y) rectangle +(0.72,0.82);
      \draw[MyNavy] (\x,\y) rectangle +(0.9,0.82);
      \node[font=\scriptsize,text=white] at (\x+0.40,\y+0.41) {80\%};
    }
  }
\end{tikzpicture}
\par\vspace{0.1cm}
Every neighborhood has the same\\racial composition as the city.
\end{column}
\begin{column}{0.49\textwidth}
\centering
\textbf{City L: small group, concentrated}\\[3pt]
100 of 1,000 residents belong to $r$.
\par\vspace{0.15cm}
\begin{tikzpicture}[x=0.88cm,y=0.75cm]
  \foreach \x in {0,...,4} {
    \foreach \y in {0,1} {
      \fill[gray!20] (\x,\y) rectangle +(0.9,0.82);
      \draw[MyNavy] (\x,\y) rectangle +(0.9,0.82);
      \node[font=\scriptsize] at (\x+0.45,\y+0.41) {0\%};
    }
  }
  \fill[MyTeal] (0,1) rectangle +(0.9,0.82);
  \draw[MyNavy] (0,1) rectangle +(0.9,0.82);
  \node[font=\scriptsize,text=white] at (0.45,1.41) {100\%};
\end{tikzpicture}
\par\vspace{0.1cm}
All members of $r$ live in one\\neighborhood, with no other group.
\end{column}
\end{columns}
\vspace{0.15cm}
{\footnotesize Each rectangle is a neighborhood. Teal area shows its share of $r$.}
\begin{block}{The comparison}
City L has complete residential separation, yet we will obtain
\alert{$S_r^{L}=0.55 < S_r^{H}=0.80$}.
\end{block}
% SAY: Residential separation and repeated same-group encounters are distinct.
% City L is completely separated at the neighborhood level, but residents can
% still encounter people in other neighborhoods under the rule on the next slide.
\end{frame}

\begin{frame}{The Same Encounter Rule in Both Cities}
\small
With probability $1/2$, draw a resident from your own neighborhood.
With probability $1/2$, draw a resident uniformly from the whole city.
\[
 P_{ij}=\tfrac12\mathbf{1}_{\{i=j\}}+\tfrac1{20},
 \qquad
 a_{ij}=P_{ij}\,\mathbb{P}_j(r),
 \qquad S_r=\rho(A).
\]
Thus $P_{ii}=0.55$ and $P_{ij}=0.05$ for $i\ne j$.
\vspace{0.1cm}
\begin{columns}[T]
\begin{column}{0.48\textwidth}
\begin{block}{City H: 80\% in every neighborhood}
Since $P$ is stochastic, $\rho(P)=1$:
\[
 A_H=0.8P,\qquad \boxed{S_r^H=0.80.}
\]
Every encounter has an 80\% chance of staying within group $r$.
\end{block}
\end{column}
\begin{column}{0.48\textwidth}
\begin{block}{City L: 100\% in neighborhood 1}
\[
 A_L=P\,\operatorname{diag}(1,0,\ldots,0).
\]
Only column 1 is nonzero, so
\[
 \boxed{S_r^L=P_{11}=0.55.}
\]
An own-group encounter requires drawing someone in neighborhood 1.
\end{block}
\end{column}
\end{columns}
\vspace{0.1cm}
{\footnotesize Fits our model: equal $D(j)$ and local weights 11 times every
nonlocal weight give exactly this $P$.}
% With the supplied exponential model, take C_ii=c_0 and
% C_ij=c_0+log(11)/beta for i!=j, with beta>0.
% Then f_ii/f_ij=11 and population-weighted row normalization gives P.
% A_L = (P e_1)e_1^T has eigenvalues P_11 and nine zeros.
% Restricting to V_r={1} gives the same nonzero eigenvalue.
\end{frame}

\begin{frame}{Why the Raw PSSI Ranking Reverses}
\small
If every neighborhood has group share $p$, then $A=pP$ and
$S_r=p$. Even uniform racial composition can yield a large raw PSSI.
\begin{center}
\renewcommand{\arraystretch}{1.15}
\begin{tabular}{lcc}
 & \textbf{City H} & \textbf{City L} \\
\hline
Citywide share of $r$ & $0.80$ & $0.10$ \\
Residential pattern & Evenly spread & Fully separated \\
\textbf{Raw PSSI} & $\mathbf{0.80}$ & $\mathbf{0.55}$ \\
Uniform-composition reference & $0.80$ & $0.10$ \\
Excess above that reference & $0.00$ & $0.45$ \\
\hline
\end{tabular}
\end{center}
\begin{block}{Motivation for normalization}
Raw PSSI reflects opportunities for repeated own-group encounters.
Comparing segregation across cities also requires accounting for
the group's population share and a specified mixing baseline.
\end{block}
{\footnotesize The uniform-composition reference is illustrative.
A random-allocation baseline $\mathbb{E}_0[\rho(A)]$ need not equal the
citywide share.}
% SAY: The raw index is answering its encounter question correctly.
% The example motivates calibration when our question is excess segregation
% relative to a city's mixing baseline. It does not prove a general ordering
% theorem, or identify population share with an expected random spectral radius.
\end{frame}
